\chapter{Mượn Đồ Học Liên Tục}
\chapterintrobi{Mang rất ít nhưng dùng rất nhiều}{Keep borrowing school things}{Đây là kiểu sống dựa vào tình nghĩa bạn bè nhiều hơn dựa vào việc chuẩn bị từ trước.}

\begin{lucbatbox}{Lục bát: Tiết nào cũng nhờ bạn}
Vào giờ em hỏi nhẹ thôi\
Bạn ơi cho mượn mấy hồi nữa nha\
Bút rồi thước, giấy rồi la\
Đồ em ít đến tưởng là gió bay\
Tình thương bạn giữ em đây\
Chứ riêng chuẩn bị chưa đầy bao nhiêu
\end{lucbatbox}

\begin{tutuyetbox}{Thất ngôn tứ tuyệt: Sống nhờ bàn bên}
Em học bằng nguồn lực của quanh em\
Mượn hết từ cây bút tới cây kèm\
Nếu tính tự lập không đi chậm thế\
Bạn bên đâu phải cứu hoài thêm
\end{tutuyetbox}

\begin{ghichu}
Mượn đồ một lần là chuyện bình thường. Mượn liên tục là biến bạn cùng bàn thành kho thiết bị lưu động.
\end{ghichu}

\begin{englishnote}
\textbf{borrow}: mượn.\par
\textbf{prepared}: chuẩn bị sẵn.\par
\textbf{Can I borrow yours again?}: Em mượn của bạn lần nữa nhé?
\end{englishnote}